\documentclass[11pt,a4paper,notitlepage]{report}
% Shared preamble for the raw chapter-4 write-ups under docs/report/insights/.
%
% Each experiment_N/experiment_N.tex is a standalone, compilable LaTeX report
% written BEFORE anything goes into final_report/chapters/chapter_4_*.tex.
% It uses the final report's fonts, units and macros so that passages, tables
% and figures lift into the chapter without reformatting.
%
%   \documentclass[11pt,a4paper]{report}
%   % Shared preamble for the raw chapter-4 write-ups under docs/report/insights/.
%
% Each experiment_N/experiment_N.tex is a standalone, compilable LaTeX report
% written BEFORE anything goes into final_report/chapters/chapter_4_*.tex.
% It uses the final report's fonts, units and macros so that passages, tables
% and figures lift into the chapter without reformatting.
%
%   \documentclass[11pt,a4paper]{report}
%   % Shared preamble for the raw chapter-4 write-ups under docs/report/insights/.
%
% Each experiment_N/experiment_N.tex is a standalone, compilable LaTeX report
% written BEFORE anything goes into final_report/chapters/chapter_4_*.tex.
% It uses the final report's fonts, units and macros so that passages, tables
% and figures lift into the chapter without reformatting.
%
%   \documentclass[11pt,a4paper]{report}
%   \input{../common/preamble}
%   \begin{document} ...
%
% Figures and tables are NOT copied: \graphicspath and \reporttable resolve into
% final_report/figures/ and final_report/tables/, so the raw write-up always
% shows the same artefacts as the chapter.  Compile with:  latexmk -pdf <file>

\usepackage[T1]{fontenc}
\usepackage[utf8]{inputenc}
\usepackage{mathptmx}
\usepackage[scaled=0.92]{helvet}
\usepackage{courier}
\usepackage[a4paper,left=30mm,right=20mm,top=25mm,bottom=25mm]{geometry}
\usepackage{microtype}
\usepackage{booktabs}
\usepackage{tabularx}
\usepackage{longtable}
\usepackage{array}
\usepackage{multirow}
\usepackage{amsmath}
\usepackage{amssymb}
\usepackage{ifthen}
\usepackage{siunitx}
\usepackage{graphicx}
\usepackage{placeins}
\usepackage{caption}
\usepackage{enumitem}
\usepackage[table]{xcolor}
\usepackage{rotating}
\PassOptionsToPackage{hyphens}{url}\usepackage{url}
\usepackage[hidelinks]{hyperref}
\usepackage[nameinlink,noabbrev]{cleveref}

\graphicspath{{../../final_report/}{../../final_report/figures/}}
% \reporttable{name} inputs final_report/tables/name.tex
\newcommand{\reporttable}[1]{\input{../../final_report/tables/#1}}

\newcommand{\repo}[1]{\texttt{\detokenize{#1}}}
\sisetup{detect-all,group-separator={,},range-phrase=--,range-units=single}
\DeclareSIUnit{\pixel}{px}
\DeclareSIUnit{\mebi}{Mi}
\captionsetup{font=small,labelfont=bf}
\setlist{nosep}
\setlength{\parskip}{4pt}
\setlength{\LTcapwidth}{\textwidth}

% ---- result-slot status (RENDERED here, unlike in the final report)
\definecolor{stcompleted}{HTML}{A5D6A7}
\definecolor{stpartial}{HTML}{FFF2CC}
\definecolor{stuntested}{HTML}{FFD8A8}
\definecolor{stblocked}{HTML}{FFCDD2}
\definecolor{stfailed}{HTML}{FFCDD2}
\definecolor{stplanned}{HTML}{E0E0E0}
\newcommand{\slotstatus}[2]{%
  \par\noindent\fcolorbox{black!40}{st#1}{\parbox{\dimexpr\textwidth-2\fboxsep-2\fboxrule}{%
    \textbf{Status: #1.}\ #2}}\par\medskip}
\newcommand{\status}[1]{\textbf{#1}}

% ---- queries: RENDERED in red here so a reader sees every open concern in place
\definecolor{flagred}{HTML}{B00000}
\newcommand{\flag}[1]{\textcolor{flagred}{#1}}
\newcommand{\moved}[1]{#1}
\newenvironment{queried}[1][Questionable]%
  {\par\addvspace{0.6ex}\begingroup\color{flagred}\noindent
   \textbf{[#1]}\ \ignorespaces}%
  {\par\endgroup\addvspace{0.6ex}}

% ---- per-block performance marks (section 4.5 diagrams)
\definecolor{perfgood}{HTML}{A5D6A7}
\definecolor{perfpoor}{HTML}{FFCDD2}
\definecolor{perfnoeffect}{HTML}{D9EAF7}
\definecolor{perfnone}{HTML}{FFF2CC}
\newcommand{\perfmark}[1]{%
  \ifthenelse{\equal{#1}{good}}{\colorbox{perfgood}{\scriptsize good}}{%
  \ifthenelse{\equal{#1}{poor}}{\colorbox{perfpoor}{\scriptsize poor}}{%
  \ifthenelse{\equal{#1}{noeffect}}{\colorbox{perfnoeffect}{\scriptsize no effect}}{%
  \colorbox{perfnone}{\scriptsize no claim}}}}}

%   \begin{document} ...
%
% Figures and tables are NOT copied: \graphicspath and \reporttable resolve into
% final_report/figures/ and final_report/tables/, so the raw write-up always
% shows the same artefacts as the chapter.  Compile with:  latexmk -pdf <file>

\usepackage[T1]{fontenc}
\usepackage[utf8]{inputenc}
\usepackage{mathptmx}
\usepackage[scaled=0.92]{helvet}
\usepackage{courier}
\usepackage[a4paper,left=30mm,right=20mm,top=25mm,bottom=25mm]{geometry}
\usepackage{microtype}
\usepackage{booktabs}
\usepackage{tabularx}
\usepackage{longtable}
\usepackage{array}
\usepackage{multirow}
\usepackage{amsmath}
\usepackage{amssymb}
\usepackage{ifthen}
\usepackage{siunitx}
\usepackage{graphicx}
\usepackage{placeins}
\usepackage{caption}
\usepackage{enumitem}
\usepackage[table]{xcolor}
\usepackage{rotating}
\PassOptionsToPackage{hyphens}{url}\usepackage{url}
\usepackage[hidelinks]{hyperref}
\usepackage[nameinlink,noabbrev]{cleveref}

\graphicspath{{../../final_report/}{../../final_report/figures/}}
% \reporttable{name} inputs final_report/tables/name.tex
\newcommand{\reporttable}[1]{\input{../../final_report/tables/#1}}

\newcommand{\repo}[1]{\texttt{\detokenize{#1}}}
\sisetup{detect-all,group-separator={,},range-phrase=--,range-units=single}
\DeclareSIUnit{\pixel}{px}
\DeclareSIUnit{\mebi}{Mi}
\captionsetup{font=small,labelfont=bf}
\setlist{nosep}
\setlength{\parskip}{4pt}
\setlength{\LTcapwidth}{\textwidth}

% ---- result-slot status (RENDERED here, unlike in the final report)
\definecolor{stcompleted}{HTML}{A5D6A7}
\definecolor{stpartial}{HTML}{FFF2CC}
\definecolor{stuntested}{HTML}{FFD8A8}
\definecolor{stblocked}{HTML}{FFCDD2}
\definecolor{stfailed}{HTML}{FFCDD2}
\definecolor{stplanned}{HTML}{E0E0E0}
\newcommand{\slotstatus}[2]{%
  \par\noindent\fcolorbox{black!40}{st#1}{\parbox{\dimexpr\textwidth-2\fboxsep-2\fboxrule}{%
    \textbf{Status: #1.}\ #2}}\par\medskip}
\newcommand{\status}[1]{\textbf{#1}}

% ---- queries: RENDERED in red here so a reader sees every open concern in place
\definecolor{flagred}{HTML}{B00000}
\newcommand{\flag}[1]{\textcolor{flagred}{#1}}
\newcommand{\moved}[1]{#1}
\newenvironment{queried}[1][Questionable]%
  {\par\addvspace{0.6ex}\begingroup\color{flagred}\noindent
   \textbf{[#1]}\ \ignorespaces}%
  {\par\endgroup\addvspace{0.6ex}}

% ---- per-block performance marks (section 4.5 diagrams)
\definecolor{perfgood}{HTML}{A5D6A7}
\definecolor{perfpoor}{HTML}{FFCDD2}
\definecolor{perfnoeffect}{HTML}{D9EAF7}
\definecolor{perfnone}{HTML}{FFF2CC}
\newcommand{\perfmark}[1]{%
  \ifthenelse{\equal{#1}{good}}{\colorbox{perfgood}{\scriptsize good}}{%
  \ifthenelse{\equal{#1}{poor}}{\colorbox{perfpoor}{\scriptsize poor}}{%
  \ifthenelse{\equal{#1}{noeffect}}{\colorbox{perfnoeffect}{\scriptsize no effect}}{%
  \colorbox{perfnone}{\scriptsize no claim}}}}}

%   \begin{document} ...
%
% Figures and tables are NOT copied: \graphicspath and \reporttable resolve into
% final_report/figures/ and final_report/tables/, so the raw write-up always
% shows the same artefacts as the chapter.  Compile with:  latexmk -pdf <file>

\usepackage[T1]{fontenc}
\usepackage[utf8]{inputenc}
\usepackage{mathptmx}
\usepackage[scaled=0.92]{helvet}
\usepackage{courier}
\usepackage[a4paper,left=30mm,right=20mm,top=25mm,bottom=25mm]{geometry}
\usepackage{microtype}
\usepackage{booktabs}
\usepackage{tabularx}
\usepackage{longtable}
\usepackage{array}
\usepackage{multirow}
\usepackage{amsmath}
\usepackage{amssymb}
\usepackage{ifthen}
\usepackage{siunitx}
\usepackage{graphicx}
\usepackage{placeins}
\usepackage{caption}
\usepackage{enumitem}
\usepackage[table]{xcolor}
\usepackage{rotating}
\PassOptionsToPackage{hyphens}{url}\usepackage{url}
\usepackage[hidelinks]{hyperref}
\usepackage[nameinlink,noabbrev]{cleveref}

\graphicspath{{../../final_report/}{../../final_report/figures/}}
% \reporttable{name} inputs final_report/tables/name.tex
\newcommand{\reporttable}[1]{\input{../../final_report/tables/#1}}

\newcommand{\repo}[1]{\texttt{\detokenize{#1}}}
\sisetup{detect-all,group-separator={,},range-phrase=--,range-units=single}
\DeclareSIUnit{\pixel}{px}
\DeclareSIUnit{\mebi}{Mi}
\captionsetup{font=small,labelfont=bf}
\setlist{nosep}
\setlength{\parskip}{4pt}
\setlength{\LTcapwidth}{\textwidth}

% ---- result-slot status (RENDERED here, unlike in the final report)
\definecolor{stcompleted}{HTML}{A5D6A7}
\definecolor{stpartial}{HTML}{FFF2CC}
\definecolor{stuntested}{HTML}{FFD8A8}
\definecolor{stblocked}{HTML}{FFCDD2}
\definecolor{stfailed}{HTML}{FFCDD2}
\definecolor{stplanned}{HTML}{E0E0E0}
\newcommand{\slotstatus}[2]{%
  \par\noindent\fcolorbox{black!40}{st#1}{\parbox{\dimexpr\textwidth-2\fboxsep-2\fboxrule}{%
    \textbf{Status: #1.}\ #2}}\par\medskip}
\newcommand{\status}[1]{\textbf{#1}}

% ---- queries: RENDERED in red here so a reader sees every open concern in place
\definecolor{flagred}{HTML}{B00000}
\newcommand{\flag}[1]{\textcolor{flagred}{#1}}
\newcommand{\moved}[1]{#1}
\newenvironment{queried}[1][Questionable]%
  {\par\addvspace{0.6ex}\begingroup\color{flagred}\noindent
   \textbf{[#1]}\ \ignorespaces}%
  {\par\endgroup\addvspace{0.6ex}}

% ---- per-block performance marks (section 4.5 diagrams)
\definecolor{perfgood}{HTML}{A5D6A7}
\definecolor{perfpoor}{HTML}{FFCDD2}
\definecolor{perfnoeffect}{HTML}{D9EAF7}
\definecolor{perfnone}{HTML}{FFF2CC}
\newcommand{\perfmark}[1]{%
  \ifthenelse{\equal{#1}{good}}{\colorbox{perfgood}{\scriptsize good}}{%
  \ifthenelse{\equal{#1}{poor}}{\colorbox{perfpoor}{\scriptsize poor}}{%
  \ifthenelse{\equal{#1}{noeffect}}{\colorbox{perfnoeffect}{\scriptsize no effect}}{%
  \colorbox{perfnone}{\scriptsize no claim}}}}}

\setcounter{chapter}{4}
\renewcommand{\thesection}{\arabic{section}}
\renewcommand{\thefigure}{I.\arabic{figure}}
\renewcommand{\thetable}{I.\arabic{table}}
% every table and figure of this document is generated by script/exp3_insights.py
\newcommand{\insdir}{../../../../output/compare/simulation/exp3-insights}
\newcommand{\ins}[1]{\input{\insdir/#1}}
\newcommand{\realdir}{../../../../output/compare/real}
\newcommand{\insreal}[1]{\input{\realdir/realv6-insights/#1}}

\title{Experiment 3: What the Pose Error Rises With\\
       \large Insight analysis for Chapter 4, Section 4.3}
\author{Bolton Athit Davies}
\date{Reporting cutoff: 22 October 2026; real-robot section added 2 November 2026}

\begin{document}
\maketitle

\noindent\textbf{Purpose of this document.} The raw write-up of Experiment~3
(\path{insights/experiment_3/}) reports what each run scored. This document asks
\emph{why}: whether the pose error rises when the robot turns, when a movable
object is in the camera's view, or after the tracker loses features, and what
the robot is doing at the moment the error starts to grow. Every question is
tested within each run, so that the run is its own control, and the answer is
stated with the evidence that supports it and the evidence that is missing.
Every number is written by \path{script/exp3_insights.py} into
\path{output/compare/simulation/exp3-insights/}; the tables below are input
from there unchanged.

\section{Questions and method}

Five questions were put to the 114 Experiment~3 runs of both testbeds (the
registry of \path{script/analyse_exp3_timelines.py}: 64 dynamic and 18 static
runs in the extreme world, 26 dynamic and 6 static in the realistic world).

\begin{enumerate}[label=Q\arabic*.]
\item Does the error grow faster while the robot turns than on straights?
\item Does it grow faster while a movable object is in the camera's view?
\item Is the onset preceded by a loss of tracked features?
\item What is the robot doing at the onset: turning or straight, object in view
      or not?
\item Is the error that these questions concern the estimator's error at all,
      or does part of it come from how the estimate is scored?
\end{enumerate}

\paragraph{Within-run contrasts.} The error series is the start-anchored planar
error of the timeline analysis: the estimate is aligned to the ground truth by a
rigid fit over the first \SIrange{10}{40}{\second}, after the initialisation
transient and before the ground truth has turned \SI{60}{\degree}, and the
planar distance to the ground truth is followed from there. For Q1 and Q2 the
samples of a run are split into two states and the growth rate of the error
(metres per second of positive and negative change) is computed in each. States
are: \emph{turning}, ground-truth yaw rate above \SI{10}{\degree\per\second};
\emph{object in view}, for the extreme bags the projected simulator truth of
Experiment~2 (image area of moving objects in the left camera), for the realistic
bags the carriers and carried loads whose poses the bag records on
\path{/world/default/pose/info}, placed within \SI{20}{\metre} and
$\pm\SI{45}{\degree}$ of the ground-truth heading with no shelf footprint on the
line of sight. Rates are taken only while the estimate is still tracking: after
the anchor window, before its first impossible-speed jump and while the error is
below \SI{3}{\metre}. A contrast needs at least \SI{3}{\second} in each state;
across runs the two rates are compared with the sign count and the Wilcoxon
signed-rank test. For Q3 the matched-feature count (ORB-SLAM3 map inliers,
VINS-Fusion tracked features) is followed around the onset as a percentage of
the run's own level before it. For Q4 the state at the onset time is read off.
The onset is the one the raw write-up and the timeline tables use: \SI{1}{\second}
of sustained excess over three times the static median error, with the ground
truth scored as the chapter scores it. The drift contrasts of Q1 and Q2 are
instead computed with the ground truth at the IMU point (Insight~1), because at
the model origin they measure the lever arm, not the estimator.

\paragraph{Lever arm (Q5).} Both estimators report the pose of the IMU: the
ORB-SLAM3 wrapper converts the camera pose with the camera-to-body extrinsic
before writing it, and VINS-Fusion writes its body state. The ground truth is the
model's odometry, which reports the model origin. In the simulated platform the
IMU sits \SI{0.338}{\metre} ahead of that origin
(\path{models/ackermann_robot/model.sdf}). No scoring step in the project
compensates this; Chapter~3 states the requirement only for the real robot's
tracker marker. The effect is tested by moving the ground truth forward along
its own heading by a lever arm $L$ and re-scoring every tracked unfiltered run
for $L$ from \SIrange{0}{0.60}{\metre}.

\FloatBarrier
\section{Insight 1: the error that appears at turns is the scoring lever arm}

Scored as the chapter scores them, the error of the tracked runs grows while the
robot turns and holds on straights; scored at the IMU point, that pattern
disappears and the whole-run ATE falls by half or more.
\Cref{tab:insight-lever,tab:insight-lever-argmin} and \cref{fig:lever} carry
the sweep.

\ins{lever_arm_sweep.tex}
\ins{lever_arm_argmin.tex}

\begin{figure}[htbp]
  \centering
  \includegraphics[width=\textwidth]{\insdir/fig_lever_sweep_light.png}
  \caption{Whole-run planar ATE of every tracked unfiltered run against the
  lever arm assumed between the model origin and the IMU. Solid lines are static
  recordings, dashed the realistic dynamic recordings. The dotted line is the
  value in the robot model.}
  \label{fig:lever}
\end{figure}

The minimum of the ATE lies at the lever arm the robot model specifies. Over the
30 tracked unfiltered runs the ATE-minimising $L^{*}$ has a median of
\SI{0.34}{\metre} and an interquartile range of \SIrange{0.32}{0.40}{\metre}
(\cref{tab:insight-lever-argmin}), for both estimators and both testbeds. The
three runs that do not follow the pattern are the ones whose error is not a
rigid offset: the two ORB-SLAM3 runs on \repo{realistic_002}, whose error is a
\SIrange{0.5}{1.0}{\metre} step taken in the first \SI{15}{\second}
(\cref{fig:traffic}), and the VINS-Fusion run on the same bag, which diverges at
its first turn. With $L=\SI{0.34}{\metre}$ the median static ATE of the extreme
world falls from \SI{0.32}{\metre} to \SI{0.10}{\metre} for VINS-Fusion and from
\SI{0.37}{\metre} to \SI{0.19}{\metre} for ORB-SLAM3; in the realistic world from
\SI{0.37}{\metre} to \SI{0.25}{\metre} and from \SI{0.32}{\metre} to
\SI{0.15}{\metre} (\cref{tab:insight-lever}).

The turning effect goes with it. At $L=0$ ORB-SLAM3's error in the realistic
world grows at a median \SI{8.7}{\milli\metre\per\second} while turning and
shrinks on straights, in eight of eight unfiltered runs; at $L=\SI{0.34}{\metre}$
the turning rate is \SI{0.0}{\milli\metre\per\second}. This is the signature of
a rigid offset between the scored point and the reference point: the offset
vector rotates with the robot, so the apparent error changes only while the
heading changes.

The sweep fits one parameter and could, on its own, absorb some estimator error.
What makes it a finding rather than a fit is that the minimum coincides with an
independently specified geometric constant, in 30 runs, across two estimators
that share no code. The residual ATE at $L^{*}$ is the estimator's.

\begin{queried}[Consequence beyond this experiment]
Every simulation ATE in Chapter~4 that aligns an estimate to the model-origin
ground truth carries this offset: the static floor of Experiment~3 and the
visual estimators of Experiment~1. The wheel and IMU baselines of Experiment~1
are referenced to the same origin and are not affected in the same way, so the
estimator ranking of that experiment has to be re-scored before it is quoted.
The correction is one line: move the ground truth to the IMU point before
alignment. It has since been applied in the scoring scripts
(\path{vio_metrics.load_gt}, 22 October 2026) and every simulation table and
figure of Experiments~1 and~3 has been regenerated; the chapter prose that
quotes the old values is pending.
\end{queried}

\FloatBarrier
\section{Insight 2: after the correction, turning does not raise the error}

With the ground truth at the IMU point, neither estimator's error grows faster
while turning than on straights, in either testbed, and the heading error does
not either. \Cref{tab:insight-contrast-extreme,tab:insight-contrast-realistic}
give the contrasts; \cref{fig:turn} plots each run.

\ins{contrast_extreme_lever0p338.tex}
\ins{contrast_realistic_lever0p338.tex}

\begin{figure}[htbp]
  \centering
  \includegraphics[width=\textwidth]{\insdir/fig_turn_contrast_lever0p338_light.png}
  \caption{Drift rate while turning against drift rate on straights, one marker
  per run, ground truth at the IMU point. Negative rates are drawn at the axis
  floor. A run above the diagonal drifts faster while turning.}
  \label{fig:turn}
\end{figure}

The static recordings of the extreme world are the only group with enough runs
for a test, and the two estimators point opposite ways there: VINS-Fusion is
faster while turning in seven of nine runs and ORB-SLAM3 in one of nine, neither
at $p<0.05$. In the realistic world the sign is mixed in every group. The
heading-error contrast is flat throughout. On these data the answer to Q1 is no.

\FloatBarrier
\section{Insight 3: in the extreme world the failure is immediate and is not
         modulated by how much moving object is in view}

Every dynamic run of the extreme world passes its onset threshold on the first
aisle, almost never while turning, and the moment-to-moment amount of moving
object in the image does not predict when ORB-SLAM3 loses tracking or when
either estimate jumps. \Cref{tab:insight-onset-extreme,tab:insight-mechanism-extreme}
and \cref{fig:epoch} carry the evidence.

\ins{onset_context_extreme.tex}
\ins{mechanism_extreme.tex}

\begin{figure}[p]
  \centering
  \includegraphics[width=\textwidth]{\insdir/fig_onset_epoch_light.png}
  \caption{What precedes the onset, at the chapter's onset definition: median
  over runs (shaded, interquartile range) of the matched-feature count relative
  to the run's own pre-onset level, the true moving-object area in view (extreme
  world) or the number of traffic agents in view (realistic world), the image
  area the filter masked, and the ground-truth yaw rate, against time from
  onset.}
  \label{fig:epoch}
\end{figure}

Moving objects are in view for \SIrange{74}{91}{\percent} of the frames of these
four bags (projected truth of Experiment~2), so a within-bag ``no object''
baseline does not exist; the comparison has to be between more and less moving
area. Split into terciles of moving area within each run, the share of frames
in which ORB-SLAM3 is not tracking OK is not higher in the top tercile than in
the bottom one (\cref{tab:insight-mechanism-extreme}), nor is the jump rate of
either estimate. The matched-feature count falls weakly with moving area
(median Spearman $\rho$ about $-0.1$ for ORB-SLAM3, in eleven of twelve runs).

The onset is on a straight. Of the 64 dynamic runs, the robot was turning at
onset in fifteen (\cref{tab:insight-onset-extreme}); the first heading change of
every route comes after the first aisle, and the error has passed threshold
before it.

The two filtered estimators reach their onset differently. For ORB-SLAM3 with
the class mask the matched features collapse to about one fifth of the run's
own level \SI{2.5}{\second} before the onset (\cref{fig:epoch}, top left): the
mask removes the features the tracker was using, and the error follows.
VINS-Fusion shows no feature loss before its onset, filtered or not; its
estimate does not drift, it breaks, with a jump rate of about 20 per second
once it has failed. What breaks it is not identified by these data; the
hypothesis that wrongly tracked moving points corrupt the visual-inertial
initialisation remains untested.

\FloatBarrier
\section{Insight 4: in the realistic world the traffic in view does not
         measurably change the drift, and the class mask costs features when
         traffic is in view}

With the carriers' true positions from the bag, the drift of the tracked
realistic runs is not consistently faster while a carrier is in the camera's
view, and ORB-SLAM3 with the class mask has fewer matched features during those
moments without a drift penalty.
\Cref{tab:insight-traffic-realistic} and \cref{fig:traffic} carry the
evidence; the per-run values are in
\path{output/compare/simulation/exp3-insights/realistic_runs_lever0p338.tex}.

\ins{traffic_realistic_lever0p338.tex}

\begin{figure}[htbp]
  \centering
  \includegraphics[width=\textwidth]{\insdir/fig_traffic_timeline_lever0p338_light.png}
  \caption{Realistic dynamic recordings, unfiltered runs: anchored planar error
  (ground truth at the IMU point) with the moments a carrier or carried load is
  in the camera's view (violet) and the moments the robot turns (grey). The
  green band is the alignment window.}
  \label{fig:traffic}
\end{figure}

Traffic is in view for a median \SI{37}{\percent} of the tracked time, so both
states are well populated; the groups of runs are not. Two unfiltered
VINS-Fusion runs and four ORB-SLAM3 runs carry the unfiltered answer, and in
them the drift is, if anything, slower while a carrier is in view (median
\SI{-0.2}{\milli\metre\per\second} against \SI{1.5}{\milli\metre\per\second}
for ORB-SLAM3, one run of four faster in view). The filtered ORB-SLAM3 group is
the only one large enough to test: there the matched-feature count is lower
while traffic is in view in eleven of sixteen runs ($p=0.05$) and the drift
rate is higher in seven of ten ($p=0.08$), with no change in the tracking-loss
share. Taken together these suggest, and do not establish, that the class mask
costs features when a carrier is in view and that the cost is paid in drift;
the unfiltered estimators show no such cost.

The one large degradation of the realistic world is not a traffic encounter.
The VINS-Fusion error on \repo{realistic_002}, \SI{1.8}{\metre} whole-run ATE,
is a single \SI{3}{\metre} step taken between \SI{29}{\second} and
\SI{33}{\second}, at the end of the first heading change and with no carrier
in the camera's view (\cref{fig:traffic}, third panel); the carrier that passed
at \SIrange{5}{12}{\second} left no trace. The same bag gives both ORB-SLAM3
runs their \SIrange{0.5}{1.0}{\metre} step at \SI{2.5}{\second} and
\SI{13}{\second}, the moments that same carrier enters and leaves the view.
Whether the step is the estimator's or the ground truth's is not settled by
these data.

\begin{queried}[Coverage of the traffic measurement]
Walking people are gz-sim actors and do not appear on the pose topic, so their
visibility is unmeasured; occlusion is tested against shelf footprints only; and
the carried loads move only while carried. The in-view state is therefore a
lower bound on the moments something moving was in the image.
\end{queried}

\FloatBarrier
\section{Insight 5: real robot (\repo{dataset_realv6_000}): ORB-SLAM3 tracks the
         figure-of-eight to \SI{0.4}{\metre}, VINS-Fusion diverges at both replay
         rates on a stale gravity norm, and one tracker sample decided the headline}

The first real recording with a usable optical-tracker reference since the
camera remount gives ORB-SLAM3 a provisional planar ATE of \SI{0.41}{\metre}
over \SI{218}{\second} at both replay rates, no tracking loss after its
initialisation resets, and two VINS-Fusion runs that diverge to kilometres from
the same moment at both replay rates, with an accelerometer bias frozen at the
difference between the configured gravity norm and the gravity this IMU
measures. \Cref{tab:realv6-cleaning,tab:realv6-orb,tab:realv6-turning,tab:realv6-vins}
and \cref{fig:realv6-gt,fig:realv6-all,fig:realv6-orb,fig:realv6-vins} carry the
evidence.

\paragraph{The recording.} \repo{dataset_realv6_000} was recorded on 31 October
2026 on the Gensurv floor: \SI{220}{\second}, 6\,601 stereo pairs, the RealSense
IMU at \SI{200}{\hertz}, wheel-speed feedback, and a figure-of-eight route of
about \SI{27}{\metre} at a median \SI{0.15}{\metre\per\second}. Four estimator
runs exist: ORB-SLAM3 and VINS-Fusion each at replay rates 1.0 and 0.5, ORB-SLAM3
on \repo{config/wil/stereo_imu.yaml} and VINS-Fusion on the 31 October
self-calibrated extrinsics (\repo{config/wil_311026_selfcal}), all without a
detector, so unfiltered. No
filtered run exists, so the filter hypothesis remains untested on the real robot;
what this slot can now say is how the unfiltered estimators behave on real
imagery against an independent reference.

\paragraph{Reporting rules.} Scoring is \path{script/eval_real_gt.py}
unchanged: time offset from the cross-correlation of the tracker's speed with the
bag's wheel speed, rigid SE(2) alignment with scale fixed, lever arm from the
tracker marker to the IMU unmeasured and set to zero, so every number is
\emph{provisional}. The tool's validity rules were fixed before the data were
read; the tracker export fails two of them (gaps above \SI{1}{\second}, a step
above \SI{0.5}{\metre}), so the tool labels every run \emph{failed (ground truth
invalid)} and the numbers below are reported as diagnostics, not as a result-slot
headline. Outputs: \path{output/compare/real/realv6*/} and
\path{output/compare/real/realv6-insights/}.

\insreal{gt_cleaning.tex}

\begin{figure}[htbp]
  \centering
  \includegraphics[width=\textwidth]{\realdir/realv6-insights/fig_gt_cleaning_light.png}
  \caption{The tracker export of \repo{dataset_realv6_000}: one sample at
  \SI{184.9}{\second} lies \SI{5.2}{\metre} off the path and back (left), which
  puts an \SI{8.7}{\metre\per\second} spike into the speed profile the time
  alignment correlates against (right).}
  \label{fig:realv6-gt}
\end{figure}

One tracker sample changed every headline. The export runs at \SI{1.7}{\hertz},
not the \SI{7}{\hertz} of the earlier recordings, with 66 gaps longer than
\SI{1}{\second}, and it contains a single sample \SI{5.2}{\metre} off the path.
That sample added \SI{10}{\metre} to the true path length (37.3 against
\SI{27.1}{\metre}), dominated the speed cross-correlation so that the time offset
came out as \SI{+4.4}{\second} with a peak correlation of 0.10, and left
ORB-SLAM3 with an ATE of \SI{1.18}{\metre} and a path ratio of 0.68
(\cref{tab:realv6-cleaning}, rows ``raw''). Removing that one sample, and
nothing else, moves the offset to \SI{-4.1}{\second} at a correlation of 0.61,
the path ratio to 0.93 and the ATE to \SI{0.41}{\metre}. The sparse rate and the
gaps remain, so the cleaned truth still fails the gap rule; the cleaning is
recorded here as a one-sample edit with its effect shown, not applied silently.

\insreal{orb_summary.tex}

\begin{figure}[htbp]
  \centering
  \includegraphics[width=0.62\textwidth]{\realdir/realv6_clean/dataset_realv6_000/trajectory_light.png}
  \caption{ORB-SLAM3 at replay rate 1.0 against the cleaned tracker truth after
  rigid SE(2) alignment. The estimate traces the figure-of-eight with its loops
  about \SI{7}{\percent} smaller than the tracker's (Sim(2) scale diagnostic
  1.09, not applied).}
  \label{fig:realv6-traj}
\end{figure}

\begin{figure}[htbp]
  \centering
  \includegraphics[width=\textwidth]{\realdir/realv6-insights/fig_all_runs_light.png}
  \caption{All four runs against the cleaned tracker truth after the same rigid
  SE(2) alignment. Left: ORB-SLAM3 at replay 1.0 and 0.5, both on the route.
  Right: VINS-Fusion at both rates, which leave the \SI{10}{\metre} floor within
  a minute; note the axis scale.}
  \label{fig:realv6-all}
\end{figure}

\begin{figure}[htbp]
  \centering
  \includegraphics[width=\textwidth]{\realdir/realv6-insights/fig_orb_timeline_light.png}
  \caption{ORB-SLAM3 planar error against the cleaned tracker truth over the
  run, both replay rates. Grey bands: the aligned estimate's heading changes
  faster than \SI{10}{\degree\per\second}. Dotted red: the frames at which
  tracking left the OK state, all inside the first \SI{13}{\second} (the
  \repo{local_mapping_bad_imu} initialisation resets).}
  \label{fig:realv6-orb}
\end{figure}

\insreal{turning.tex}

ORB-SLAM3 tracks the whole route at both replay rates. After two (rate 1.0) or
three (rate 0.5) map resets during the first \SI{13}{\second}, while the
inertial initialisation settles, tracking never leaves the OK state again; the
error is a slow wave between \SI{0.05}{} and \SI{0.85}{\metre} with its crest
around \SIrange{120}{140}{\second} and its trough at \SI{195}{\second}, not a
drift that accumulates (\cref{fig:realv6-orb}). Turning does not raise it: the
error shrinks slightly during the heading changes in both runs
(\cref{tab:realv6-turning}), the same answer the simulation gave once the lever
arm was removed. Real time costs frames but not accuracy here: at rate 1.0 the
node dropped 3\,390 of 6\,476 frames and scored \SI{0.41}{\metre}; at rate 0.5
it dropped 619 and scored \SI{0.57}{\metre}, with the same shape of error curve.

What the \SI{0.4}{\metre} is made of is not settled. The loops of the estimate
are a uniform \SI{7}{\percent} smaller than the tracker's
(\cref{fig:realv6-traj}; Sim(2) scale 1.09; path ratio 0.93), which a
stereo-baseline error, a tracker scale error or the unmeasured lever arm could
each produce; the least-squares lever-arm fit (\SI{0.31}{\metre} behind the IMU,
not applied) would bring the ATE to \SI{0.31}{\metre}, which is the size of the
effect the unmeasured arm can hide.

\insreal{vins_summary.tex}

\begin{figure}[htbp]
  \centering
  \includegraphics[width=\textwidth]{\realdir/realv6-insights/fig_vins_divergence_light.png}
  \caption{VINS-Fusion on \repo{dataset_realv6_000} at both replay rates:
  distance of the estimate from its origin (log scale), estimated speed, and the
  accelerometer-bias magnitude the estimator holds. The dashed line is the
  difference between what the IMU measures at rest in this bag
  (\SI{9.878}{\metre\per\second\squared}) and the gravity norm the
  configuration fixes (\num{9.2642}).}
  \label{fig:realv6-vins}
\end{figure}

VINS-Fusion diverges at both replay rates from the same moment, so the frame
rate is not what breaks it. The two nodes received 1\,047 and 2\,264 of the
bag's 6\,601 frames (one in six and one in three), initialised at
\SI{0.33}{\second} and never raised a reset; in both the estimated speed passes
\SI{1}{\metre\per\second} at \SIrange{37}{38}{\second}, while the robot never
exceeds \SI{0.9}{\metre\per\second}, and from there the position runs away to
about \SI{5}{\kilo\metre} (\cref{tab:realv6-vins}, \cref{fig:realv6-vins}).
The quantity that identifies the cause is the accelerometer bias: both runs
settle it at \SIrange{0.68}{0.72}{\metre\per\second\squared} and then hold it,
which is within \SI{0.1}{\metre\per\second\squared} of the gap between the
gravity norm the configuration fixes, \num{9.2642}, measured on an August
recording of an earlier IMU, and the \SI{9.878}{\metre\per\second\squared}
this bag's IMU reads at rest. The estimator absorbs the mismatch into the bias
while the visual constraints are strong, and once they thin out the mis-scaled
gravity integrates into velocity. ORB-SLAM3 on the same bag is unaffected
because its launch rescales the accelerometer (\repo{imu_accel_scale}) instead
of fixing a gravity norm. The frame starvation is real and secondary: it
weakens the visual constraint that was holding the bias in place, and it sits on
VINS-Fusion's input path, since the ORB-SLAM3 node received 6\,476 frames from
the same bag.

\begin{queried}[What this slot still needs before it is a result]
A tracker export at its nominal \SI{7}{\hertz} without gaps, or an explicit rule
for scoring on a sparse reference; the marker-to-IMU lever arm measured, which
bounds \SI{0.1}{\metre} of the ORB-SLAM3 figure; a VINS-Fusion run with
\repo{g_norm} set to what this IMU measures (or the accelerometer rescaled as
ORB-SLAM3 does) and with the full frame stream delivered; and at least one
filtered run per estimator, without which Experiment~3's hypothesis has no
real-robot test.
\end{queried}

\FloatBarrier
\section{What these data cannot answer}

\begin{itemize}
\item \textbf{Whether the filter helps.} No admissible filtered-versus-unfiltered
      pair exists in either simulated testbed, and no filtered run exists on the
      real robot; nothing here changes that.
\item \textbf{Why VINS-Fusion fails in the extreme world.} The failure is
      immediate and not preceded by feature loss or by a rise in moving-object
      area; the mechanism is not identified.
\item \textbf{The traffic effect with confidence.} Two to four unfiltered runs
      per estimator in the realistic world, with the people unmeasured, bound
      the effect rather than estimate it.
\item \textbf{The feature-collapse signal, robustly.} The onset threshold is
      three times the static median error, and that median carries the lever
      arm. Recomputing the thresholds at the IMU point lowers them, moves the
      onsets earlier, and weakens the ORB-SLAM3+YOLO feature collapse of
      \cref{fig:epoch} to a dip to about half in four runs
      (\path{fig_onset_epoch_lever0p338_light.png}). The collapse is real at
      the chapter's definition and threshold-dependent in size.
\end{itemize}

\section{Summary}

\begin{table}[htbp]
\centering\small
\caption{Questions, answers and the strength of the evidence.}
\label{tab:insight-summary}
\begin{tabularx}{\textwidth}{lXXl}
\toprule
 & Answer & Evidence & Confidence \\
\midrule
Q5 lever arm & Yes: \SI{0.34}{\metre} of forward offset is in every simulation ATE & minimum of the sweep at the SDF value in 30 runs, two estimators, two testbeds & high \\
Q1 turning & No, once the lever arm is removed & contrasts of \cref{tab:insight-contrast-extreme,tab:insight-contrast-realistic}; heading error flat & medium (small groups) \\
Q2 object in view, extreme & Not measurable within a bag; moving area does not modulate tracking loss or jumps & \cref{tab:insight-mechanism-extreme} & medium \\
Q2 traffic in view, realistic & No effect on the unfiltered drift; the class mask loses features and gains drift while traffic is in view (suggestive) & \cref{tab:insight-traffic-realistic}, \cref{fig:traffic} & low (n) \\
Q3 features before onset & ORB-SLAM3+YOLO yes (collapse to one fifth); VINS-Fusion no & \cref{fig:epoch} & medium \\
Q4 state at onset & On the first straight, not turning; objects in view & \cref{tab:insight-onset-extreme} & high \\
Real robot, ORB & Tracks the whole \SI{218}{\second} route; provisional ATE \SI{0.41}{\metre}; no turning effect; \SI{7}{\percent} scale shortfall unexplained & \cref{tab:realv6-orb,fig:realv6-orb} & medium (provisional, sparse truth) \\
Real robot, VINS & Diverges at both rates from \SI{37}{\second}; accelerometer bias frozen at the \num{0.6} gap between configured gravity norm and the IMU's; frame starvation secondary & \cref{tab:realv6-vins}, \cref{fig:realv6-vins} & high (mechanism), medium (attribution) \\
Real robot, truth & One tracker sample set the raw ATE to \SI{1.18}{\metre}; truth at \SI{1.7}{\hertz} fails the gap rule & \cref{tab:realv6-cleaning} & high \\
\bottomrule
\end{tabularx}
\end{table}

\paragraph{Consequences for Chapter 4.} Re-score every simulation ATE with the
ground truth moved to the IMU point before it is quoted, in Experiments~1 and~3.
Report the extreme world as an immediate-failure testbed in which the filter
cannot be evaluated, not as a drift testbed. Report the realistic world's
traffic effect as bounded, not measured, and the class mask's feature cost as
the one mechanism these data show. For the real-robot slot, quote ORB-SLAM3's
\SI{0.41}{\metre} as provisional with the sparse-truth and lever-arm caveats
attached, report the VINS-Fusion failure as a configuration defect (stale gravity
norm) compounded by frame delivery rather than an estimator result, and keep the slot \emph{partial}
until a filtered run and a full-rate truth exist.

\end{document}
